\section{Design a Vending Machine}
\label{sec:case-vending}

\problemstatement{Design a vending machine that holds products in slots,
accepts coins/notes, dispenses a selected product, returns change, and
handles out-of-stock and insufficient-funds cases.}

\begin{patternbox}[Clarify Before Designing]
Ask: What payment types (cash, card)? Should it make exact change, and
what if it cannot? Can a selection be cancelled mid-transaction? For this
design assume: cash only, refunds allowed, and a clear lifecycle from idle
to dispensing. The lifecycle is the heart of the problem.
\end{patternbox}

\subsection*{Identify the entities}

\texttt{VendingMachine} (the context/state machine), \texttt{Product},
\texttt{Slot}/\texttt{Inventory} (stock levels), \texttt{Coin}/payment
handling, and a set of \texttt{State} classes -- \texttt{IdleState},
\texttt{HasMoneyState}, \texttt{DispensingState}. The machine's behaviour
for the same action (``insert coin,'' ``select'') differs entirely by
state, which is the signal for the State pattern.

\begin{ideabox}[State Pattern Is the Backbone]
Model the machine as a \textbf{State} machine: each state handles
\texttt{insertCoin}, \texttt{selectProduct}, and \texttt{dispense}
appropriately and transitions the machine forward. This replaces a giant
\texttt{switch(currentState)} in every method and makes illegal
transitions (dispensing before paying) structurally impossible.
\end{ideabox}

\subsection*{Class design (C++ skeleton)}

\begin{lstlisting}
#include <bits/stdc++.h>
using namespace std;

class VendingMachine;
struct State {
    virtual void insertCoin(VendingMachine&, int amount) = 0;
    virtual void selectProduct(VendingMachine&, const string& code) = 0;
    virtual ~State() = default;
};

class VendingMachine {
    unique_ptr<State> state;
public:
    int balance = 0;
    unordered_map<string,int> stock;            // code -> quantity
    unordered_map<string,int> price;            // code -> price

    VendingMachine();                           // starts in IdleState
    void setState(unique_ptr<State> s) { state = move(s); }
    void insertCoin(int a) { state->insertCoin(*this, a); }
    void selectProduct(const string& c) { state->selectProduct(*this, c); }
};

struct IdleState : State {
    void insertCoin(VendingMachine& m, int a) override;
    void selectProduct(VendingMachine&, const string&) override {
        cout << "Insert money first\n";
    }
};
struct HasMoneyState : State {
    void insertCoin(VendingMachine& m, int a) override { m.balance += a; }
    void selectProduct(VendingMachine& m, const string& code) override {
        if (m.stock[code] == 0) { cout << "Out of stock\n"; return; }
        if (m.balance < m.price[code]) { cout << "Insufficient funds\n"; return; }
        m.stock[code]--;
        int change = m.balance - m.price[code];
        m.balance = 0;
        cout << "Dispensed " << code << ", change " << change << "\n";
        m.setState(make_unique<IdleState>());   // back to idle
    }
};

void IdleState::insertCoin(VendingMachine& m, int a) {
    m.balance += a;
    m.setState(make_unique<HasMoneyState>());   // transition
}
VendingMachine::VendingMachine() : state(make_unique<IdleState>()) {}

int main() {
    VendingMachine m;
    m.stock["A1"] = 2; m.price["A1"] = 30;
    m.selectProduct("A1");                      // "Insert money first"
    m.insertCoin(20); m.insertCoin(20);
    m.selectProduct("A1");                      // dispenses, change 10
    return 0;
}
\end{lstlisting}

\begin{takeawaybox}
The vending machine is the textbook \textbf{State} pattern problem:
recognise that the same operations behave differently per mode and give
each mode a class that acts and transitions. Follow-ups: card payments
(Strategy), exact-change detection, and cancel/refund as another
transition.
\end{takeawaybox}
